\documentclass[11pt,a4paper]{article}

\usepackage[T1]{fontenc}
\usepackage[utf8]{inputenc}
\usepackage[english]{babel}
\usepackage{geometry}
\usepackage{listings}
\usepackage{xcolor}
\usepackage{booktabs}
\usepackage{longtable}
\usepackage{array}
\usepackage{hyperref}
\usepackage{fancyhdr}
\usepackage{titlesec}
\usepackage{enumitem}
\usepackage{graphicx}
\usepackage{amsmath}
\usepackage{parskip}

\geometry{
    top=25mm, bottom=25mm,
    left=25mm, right=25mm
}

% ---------------------------------------------------------------------------
% Colours
% ---------------------------------------------------------------------------
\definecolor{codebg}{RGB}{245,245,245}
\definecolor{codeframe}{RGB}{200,200,200}
\definecolor{kwcolor}{RGB}{0,0,180}
\definecolor{commentcolor}{RGB}{80,120,80}
\definecolor{stringcolor}{RGB}{160,32,32}
\definecolor{shellprompt}{RGB}{100,100,100}

% ---------------------------------------------------------------------------
% Listings style
% ---------------------------------------------------------------------------
\lstdefinestyle{shell}{
    backgroundcolor=\color{codebg},
    frame=single,
    framerule=0.4pt,
    rulecolor=\color{codeframe},
    basicstyle=\ttfamily\small,
    breaklines=true,
    breakatwhitespace=false,
    showstringspaces=false,
    tabsize=4,
    captionpos=b,
    xleftmargin=6pt,
    xrightmargin=6pt,
    aboveskip=8pt,
    belowskip=8pt,
}

\lstdefinestyle{vhdl}{
    backgroundcolor=\color{codebg},
    frame=single,
    framerule=0.4pt,
    rulecolor=\color{codeframe},
    basicstyle=\ttfamily\small,
    keywordstyle=\color{kwcolor}\bfseries,
    commentstyle=\color{commentcolor}\itshape,
    stringstyle=\color{stringcolor},
    morekeywords={library,use,entity,port,architecture,signal,process,
                  begin,end,if,then,else,elsif,rising_edge,std_logic,
                  std_logic_vector,unsigned,integer,constant,when,
                  inout,in,out,downto,others,not,and,or},
    breaklines=true,
    showstringspaces=false,
    tabsize=4,
    captionpos=b,
    xleftmargin=6pt,
    xrightmargin=6pt,
    aboveskip=8pt,
    belowskip=8pt,
}

\lstset{style=shell}

% ---------------------------------------------------------------------------
% Page style
% ---------------------------------------------------------------------------
\pagestyle{fancy}
\fancyhf{}
\fancyhead[L]{\small VERA RBL-C64 -- FPGA Build \& Programming Guide}
\fancyhead[R]{\small Rev~6}
\fancyfoot[C]{\thepage}
\renewcommand{\headrulewidth}{0.4pt}

% ---------------------------------------------------------------------------
% Section formatting
% ---------------------------------------------------------------------------
\titleformat{\section}{\large\bfseries}{}{0em}{\thesection\quad}
\titleformat{\subsection}{\normalsize\bfseries}{}{0em}{\thesubsection\quad}

% ---------------------------------------------------------------------------
% Convenience macros
% ---------------------------------------------------------------------------
\newcommand{\cmd}[1]{\texttt{#1}}
\newcommand{\sig}[1]{\texttt{#1}}
\newcommand{\file}[1]{\texttt{#1}}
\newcommand{\hex}[1]{\texttt{\$#1}}

% ---------------------------------------------------------------------------
\begin{document}
% ---------------------------------------------------------------------------

\begin{titlepage}
    \centering
    \vspace*{30mm}
    {\Huge\bfseries VERA RBL-C64\\[4mm]
     FPGA Build \& Programming Guide\par}
    \vspace{10mm}
    {\large iCE40UP5K -- Rev~6\par}
    \vspace{8mm}
    {\normalsize
        Commodore 64 expansion-port interface:\\
        VERA registers in IO2 and GeoRAM 256\,KB on external SRAM\par}
    \vfill
    {\small Last updated: \today\par}
\end{titlepage}

\tableofcontents
\newpage

% ===========================================================================
\section{Overview}
% ===========================================================================

This document describes the complete workflow for building and programming
the FPGA firmware for the VERA RBL-C64 module.  The firmware implements the
Commodore 64 expansion-port bus interface on a Lattice iCE40UP5K\@.  It drives
the VERA chip select for the registers in \sig{IO2} and implements a
\textbf{GeoRAM 256\,KB} on an external AS6C4008 SRAM, with the GeoRAM bank
registers inside the FPGA.

There is a single bitstream (\file{C6502\_Interface.bin}); no generics or
build variants have to be chosen.

\subsection{Architecture summary}

\begin{itemize}[noitemsep]
    \item \textbf{Qualification by \sig{/IO1} and \sig{/IO2}}: every access
          is qualified by the C64 \sig{/IO1} (port pin~7, \hex{DE00}--\hex{DEFF})
          and \sig{/IO2} (port pin~10, \hex{DF00}--\hex{DFFF}) lines.  The C64
          PLA asserts them only for CPU accesses (\sig{AEC} high), so VIC-II
          fetches, whose addresses also appear on the port, never reach VERA
          or the SRAM.  \sig{/IOx} already decodes A15--A8, so only
          \textbf{A7--A0} are wired to the FPGA.
    \item \textbf{VERA registers} (\hex{DF00}--\hex{DF1F}, 32\,B): the FPGA
          asserts \sig{VERA\_CS} for one bus cycle only when \sig{/IO2} is
          low, A7--A5 are zero, the VERA select latch at \hex{DFFD} is set
          and PHI2 is high.  The VERA FPGA drives the data bus.
    \item \textbf{VERA select latch} (\hex{DFFD}): a write of any non-zero
          value selects the VERA, \texttt{\$00} deselects it.  The latch
          always responds and selects only the VERA.
    \item \textbf{GeoRAM} (\hex{DE00}--\hex{DEFF} window, 256\,KB total):
          served by an external 512K$\times$8 SRAM (AS6C4008-55).  Window and
          page/block registers (\hex{DFFE}/\hex{DFFF}) are \textbf{always
          active}, independent of the \hex{DFFD} latch, like the original
          cartridge: existing GeoRAM software works without touching
          \hex{DFFD}.
    \item \textbf{PHI2 synchronisation}: PHI2 is treated as a data input
          and synchronised to the 25\,MHz FPGA clock via a two-stage
          flip-flop chain, producing a single-cycle \cmd{phi2\_rise}
          pulse.  Register writes are sampled \cmd{SAMPLE\_CLKS} clocks later
          (Section~\ref{sec:timing}).
    \item \textbf{No BRAM}: the design uses 0 of the 30 EBR blocks.
\end{itemize}

\subsection{Memory map}

\begin{center}
\begin{tabular}{llp{5.2cm}l}
\toprule
\textbf{Address} & \textbf{Size} & \textbf{Function} & \textbf{Implementation} \\
\midrule
\hex{DE00}--\hex{DEFF} & 256\,B & GeoRAM data window (one page),
                                   always active & External SRAM \\
\hex{DF00}--\hex{DF1F} & 32\,B  & VERA registers, \sig{VERA\_CS}
                                   when the VERA is selected & VERA FPGA \\
\hex{DFFD}             & 1\,B   & VERA select latch (write only):
                                   $\neq$\,\texttt{\$00} selects,
                                   \texttt{\$00} deselects & Register \\
\hex{DFFE}             & 1\,B   & GeoRAM page register (write only,
                                   0--63, D5--D0) & Register \\
\hex{DFFF}             & 1\,B   & GeoRAM block register (write only,
                                   0--15, D3--D0) & Register \\
\bottomrule
\end{tabular}
\end{center}

\hex{DF20}--\hex{DFFC} are not decoded.  The GeoRAM registers are decoded
exactly at \hex{DFFE}/\hex{DFFF}, without the \hex{DF80}--\hex{DFFD}
mirrors of some GeoRAM cartridges, so the latch at \hex{DFFD} does not
collide with them.  Latch and GeoRAM registers power up at zero (VERA
deselected) and are not cleared by a C64 reset, like the original GeoRAM.

\begin{lstlisting}[style=shell,caption={Typical use from 6510 code}]
LDA #$04 : STA $DFFF   ; GeoRAM block 4 (always active)
LDA #$23 : STA $DFFE   ; GeoRAM page 35
LDA $DE45              ; GeoRAM byte $12345
LDA #$01 : STA $DFFD   ; select the VERA at $DF00-$DF1F
LDA #$00 : STA $DFFD   ; deselect the VERA
\end{lstlisting}

\subsection{VERA chip select}

\begin{lstlisting}[style=vhdl,caption={VERA decode in \file{C6502\_Interface.vhd}}]
sel_vera <= '1' when IO2_N='0' and A(7 downto 5) = "000" and vera_sel='1' else '0';
VERA_CS  <= '0' when sel_vera='1' and PHI2='1' else '1';
\end{lstlisting}

\sig{VERA\_CS} is never asserted for \hex{DF20}--\hex{DFFF} (latch and
GeoRAM registers included), for \hex{DExx}, for any other address, for
VIC-II cycles (no \sig{/IO2}), or for \hex{DF00}--\hex{DF1F} while
\hex{DFFD} holds zero.  PHI2 gating hides the address decode glitches of
PHI1.

% ===========================================================================
\section{Prerequisites}
% ===========================================================================

\subsection{FPGA toolchain}

While some tools are available via \cmd{apt} on Debian/Ubuntu, they are often outdated (especially on Debian 11). The recommended way to set up the toolchain is using the \textbf{OSS CAD Suite}, which provides pre-compiled, up-to-date versions of Yosys, nextpnr, GHDL, and all required plugins.

\subsubsection{Installing OSS CAD Suite (Recommended for Debian 11)}

\begin{enumerate}
    \item Download the latest Linux x64 archive from: \\
          \url{https://github.com/YosysHQ/oss-cad-suite-build/releases/latest}
    \item Extract the archive (e.g., to \file{\$HOME/oss-cad-suite}).
    \item Add the \file{bin} directory to your \cmd{PATH}:
\end{enumerate}

\begin{lstlisting}[style=shell]
export PATH="$HOME/oss-cad-suite/bin:$PATH"
\end{lstlisting}

This suite includes \cmd{yosys}, \cmd{nextpnr-ice40}, \cmd{icepack}, \cmd{iceprog}, and \cmd{ghdl} with the \cmd{ghdl-yosys-plugin} already integrated.

\subsubsection{Alternative: Manual installation via apt (Debian 12+)}

On newer distributions, you can try:

\begin{lstlisting}[style=shell]
sudo apt install fpga-icestorm nextpnr-ice40 yosys ghdl
\end{lstlisting}

\textbf{Note:} if installing via \cmd{apt}, you will likely need to build the \cmd{ghdl-yosys-plugin} from source as described below, because it is rarely packaged.

\begin{lstlisting}[style=shell]
git clone https://github.com/ghdl/ghdl-yosys-plugin
cd ghdl-yosys-plugin
make
sudo make install
\end{lstlisting}

\subsection{Flash programming tools}

On modern distributions (e.g. Debian 12+), \cmd{openFPGALoader} can be installed via package manager:

\begin{lstlisting}[style=shell]
sudo apt install openFPGALoader
\end{lstlisting}

On older distributions like \textbf{Debian 11 (Bullseye)}, the package is not available and must be built from source.

\subsubsection{Building openFPGALoader from source}

First, install the required dependencies and build tools:

\begin{lstlisting}[style=shell]
sudo apt install git cmake make g++ pkg-config libusb-1.0-0-dev \
                 libhidapi-dev libgpiod-dev libusbmuxd-dev
\end{lstlisting}

Then, clone the repository and build the project:

\begin{lstlisting}[style=shell]
git clone https://github.com/trabucayre/openfpgaloader
cd openfpgaloader
mkdir build
cd build
# Use -DCMAKE_INSTALL_PREFIX=$HOME to install in your home directory
# or omit it to install system-wide (requires sudo make install)
cmake .. -DCMAKE_INSTALL_PREFIX=$HOME
make -j$(nproc)
make install
\end{lstlisting}

Ensure \file{\$HOME/bin} is in your \cmd{PATH} if you used the home directory prefix.

\subsubsection{Iceprog}
The \cmd{iceprog} tool is already included in the \cmd{fpga-icestorm} package and supports FTDI-based programmers (e.g. FT2232H).

\subsection{Tool roles at a glance}

\begin{center}
\begin{tabular}{llll}
\toprule
\textbf{Tool} & \textbf{Input} & \textbf{Output} & \textbf{Purpose} \\
\midrule
\cmd{ghdl} (plugin)         & VHDL sources         & RTLIL            & VHDL elaboration \\
\cmd{yosys}                 & RTLIL                & \file{.json}     & Synthesis \\
\cmd{nextpnr-ice40}         & \file{.json + .pcf}  & \file{.asc}      & Place \& route \\
\cmd{icepack}               & \file{.asc}          & \file{.bin}      & Bitstream packing \\
\cmd{openFPGALoader} / \cmd{iceprog} & \file{.bin} & SPI flash        & Programming \\
\bottomrule
\end{tabular}
\end{center}

% ===========================================================================
\section{Build Flow}
% ===========================================================================

The design is a single VHDL file with no ROM to pre-load, so the build
starts directly from synthesis.

\subsection{Step 1 -- Synthesis}

\begin{lstlisting}[style=shell,caption={Synthesis with Yosys and the GHDL plugin}]
cd firmware/FPGA
yosys -m ghdl -p \
    "ghdl --std=08 C6502_Interface.vhd -e C6502_Interface; \
     synth_ice40 -top C6502_Interface -json C6502_Interface.json"
\end{lstlisting}

The \cmd{-m ghdl} flag loads the GHDL plugin into Yosys, which handles
VHDL analysis and elaboration.  The entity has one generic,
\cmd{SAMPLE\_CLKS} (default 6), that sets the register write sample point
(Section~\ref{sec:timing}); the default needs no override.

Last measured results (OSS CAD Suite 2025-10-08):

\begin{center}
\begin{tabular}{lrr}
\toprule
\textbf{Resource} & \textbf{Used} & \textbf{Available} \\
\midrule
ICESTORM\_LC          &  40 & 5280 \\
ICESTORM\_RAM (EBR)   &   0 &   30 \\
SB\_IO (user I/O)     &  35 &   39 \\
ICESTORM\_SPRAM       &   0 &    4 \\
\midrule
\textbf{Fmax (CLK)}   & \multicolumn{2}{l}{\textbf{89.54\,MHz} post-route (target 25\,MHz)} \\
\bottomrule
\end{tabular}
\end{center}

These figures predate the change that made the GeoRAM independent of the
\hex{DFFD} latch; the logic is slightly smaller now.  Re-measure after the
next build.

\subsection{Step 2 -- Place \& Route}

Place \& route assigns logic cells and routes signals to physical
iCE40 resources, guided by the pin constraints file.

\begin{lstlisting}[style=shell,caption={Place \& route with nextpnr-ice40}]
nextpnr-ice40 --up5k --package sg48 \
    --json C6502_Interface.json \
    --pcf board-pin-mapping.pcf \
    --asc C6502_Interface.asc
\end{lstlisting}

\textbf{Important:} all pin numbers in \file{board-pin-mapping.pcf} are
placeholders.  Edit the file to match your actual schematic before this
step (see Section~\ref{sec:pins}).

\subsection{Step 3 -- Pack the bitstream}

\cmd{icepack} converts the ASCII intermediate format produced by
nextpnr into the binary bitstream that is written to the SPI flash.

\begin{lstlisting}[style=shell,caption={Pack the binary bitstream}]
icepack C6502_Interface.asc C6502_Interface.bin
\end{lstlisting}

The \file{.bin} file is what gets programmed into the SPI flash.
Bitstream size is device-dependent (not design-dependent) on the iCE40:
\textbf{$\approx$104\,KB} for the iCE40UP5K SG48.

\subsection{Step 4 -- Program the SPI flash}

The iCE40UP5K loads its configuration from an external SPI NOR flash
at power-up.  The programmer writes the bitstream into that flash.

\subsubsection{Using openFPGALoader (recommended)}

\begin{lstlisting}[style=shell,caption={Program with openFPGALoader}]
openFPGALoader -b ice40_generic C6502_Interface.bin
\end{lstlisting}

Select the programmer cable explicitly if needed:

\begin{lstlisting}[style=shell]
openFPGALoader --cable ft2232      C6502_Interface.bin   # FTDI FT2232H
openFPGALoader --cable ft232       C6502_Interface.bin   # FTDI FT232H
openFPGALoader --cable tigard      C6502_Interface.bin   # Tigard
openFPGALoader --cable raspberrypi C6502_Interface.bin   # Raspberry Pi GPIO
\end{lstlisting}

\subsubsection{Using iceprog (FTDI FT2232H)}

\begin{lstlisting}[style=shell,caption={Program with iceprog}]
iceprog C6502_Interface.bin

# Erase + write (useful if previous write was partial)
iceprog -e 1048576 C6502_Interface.bin
\end{lstlisting}

\subsection{One-command build and flash}

The \file{firmware/FPGA/Makefile} wraps all steps:

\begin{lstlisting}[style=shell]
cd firmware/FPGA
make                   # synthesise + P&R + pack -> C6502_Interface.bin
make flash             # write the bitstream via openFPGALoader
make flash-iceprog     # write the bitstream via iceprog (FTDI FT2232H)
make doc               # regenerate this PDF
make clean             # remove all generated files
\end{lstlisting}

% ===========================================================================
\section{Programmer Wiring}
% ===========================================================================

The SPI flash is shared between the FPGA (configuration bus at boot)
and the programmer (writes the bitstream).  During programming, the
FPGA is held in reset by asserting \sig{CRESET\_B} low.

\begin{center}
\begin{tabular}{lll}
\toprule
\textbf{Programmer signal} & \textbf{iCE40UP5K pin} & \textbf{SPI flash pin} \\
\midrule
MOSI  & SPI\_SI   & DI  \\
MISO  & SPI\_SO   & DO  \\
SCK   & SPI\_SCK  & CLK \\
CS\#  & --        & CS\# \\
CRST\# & CRESET\_B & -- \\
CDONE & CDONE     & -- \\
\bottomrule
\end{tabular}
\end{center}

Recommended SPI flash: Winbond W25Q16 (2\,MB) or W25Q32 (4\,MB).

Supported programmers:

\begin{center}
\begin{tabular}{lll}
\toprule
\textbf{Hardware} & \textbf{Interface} & \textbf{Tool} \\
\midrule
FTDI FT232H breakout   & SPI bitbang & \cmd{iceprog}, \cmd{openFPGALoader} \\
FTDI FT2232H breakout  & MPSSE SPI   & \cmd{iceprog}, \cmd{openFPGALoader} \\
Tigard                 & JTAG/SPI    & \cmd{openFPGALoader} \\
Raspberry Pi (GPIO)    & SPI bitbang & \cmd{openFPGALoader} \\
dirtyJTAG (STM32F103)  & JTAG        & \cmd{openFPGALoader} \\
\bottomrule
\end{tabular}
\end{center}

% ===========================================================================
\section{Reset Sequencing}
\label{sec:reset}
% ===========================================================================

The VERA RBL-C64 module requires a defined power-on boot sequence: the
Commodore 64 and the VERA graphics chip must remain in reset until all
FPGA configurations have completed.  This is implemented with a single
74LVC1G08 logic gate --- no MCU or firmware is needed.

\subsection{Boot sequence}

\begin{enumerate}[noitemsep]
    \item \textbf{Power-on:} both \sig{VERA\_RESET\#} and
          \sig{C64\_RESET\#} are held low.
    \item \textbf{FPGA ready} (\sig{FPGA\_CDONE} goes high,
          $\approx$50--100\,ms after power-on): \sig{VERA\_RESET\#} is
          released, allowing VERA to start loading its own bitstream
          from its SPI flash.
    \item \textbf{VERA ready} (\sig{VERA\_CDONE} goes high,
          $\approx$200--500\,ms later): \sig{C64\_RESET\#} is released
          and the C64 boots normally.
\end{enumerate}

The natural SPI flash loading times provide all necessary delays; no RC
timer or additional delay circuit is required.

\subsection{Logic}

\begin{center}
\begin{tabular}{ll}
\toprule
\textbf{Signal} & \textbf{Boolean expression} \\
\midrule
\sig{VERA\_RESET\#} & \sig{FPGA\_CDONE} \\
\sig{C64\_RESET\#}  & \sig{FPGA\_CDONE}~AND~\sig{VERA\_CDONE} \\
\bottomrule
\end{tabular}
\end{center}

\begin{center}
\begin{tabular}{ccccl}
\toprule
\sig{FPGA\_CDONE} & \sig{VERA\_CDONE} &
\sig{VERA\_RST\#} & \sig{C64\_RST\#} & \textbf{Phase} \\
\midrule
0 & 0 & 0 & 0 & Power-on / FPGA loading \\
1 & 0 & 1 & 0 & FPGA ready, VERA loading \\
1 & 1 & 1 & 1 & System ready \\
\bottomrule
\end{tabular}
\end{center}

\subsection{Circuit -- 74LVC1G08 (SOT-23-5)}

A single \textbf{74LVC1G08} (2-input AND gate, SOT-23-5) implements
\sig{C64\_RESET\#}.  \sig{VERA\_RESET\#} connects directly from
\sig{FPGA\_CDONE} through a 33\,$\Omega$ series resistor.

\begin{lstlisting}[style=shell]
VCC --[4k7]--+-------- FPGA_CDONE (open-drain)
             |
             +--[33R]- VERA_RESET#
             |
             +-- A -\
                     +--- 74LVC1G08 ---[33R]--- C64_RESET#
             +-- B -/
             |
VCC --[4k7]--+-------- VERA_CDONE (open-drain)
\end{lstlisting}

Both CDONE outputs are open-drain; they require an external pullup to
VCC\@.  The 33\,$\Omega$ series resistors limit current if an external
driver (e.g.\ a reset button on the C64 \sig{/RESET} line, which pulls it
to GND) conflicts with the 74LVC output.

\textbf{Why 74LVC and not 74HC?}  The 74LVC family has \emph{5\,V
tolerant inputs} while operating at 3.3\,V VCC\@.  The CDONE signals
are at 3.3\,V; the C64 \sig{/RESET} line belongs to a 5\,V TTL system.
The 74LVC output at 3.3\,V is correctly recognised as HIGH by 5\,V TTL
logic (V$_{\text{IH}} = 2.0$\,V).  The 74HC family is not suitable for
this mixed-voltage interface.

\subsection{Bill of materials}

\begin{center}
\begin{tabular}{lllp{4.5cm}}
\toprule
\textbf{Ref} & \textbf{Value} & \textbf{Package} & \textbf{Function} \\
\midrule
U1  & 74LVC1G08      & SOT-23-5 & 2-input AND gate \\
R1  & 4.7\,k$\Omega$ & 0402     & Pullup for FPGA\_CDONE \\
R2  & 4.7\,k$\Omega$ & 0402     & Pullup for VERA\_CDONE \\
R3  & 33\,$\Omega$   & 0402     & Series resistor on VERA\_RESET\# \\
R4  & 33\,$\Omega$   & 0402     & Series resistor on C64\_RESET\# \\
\bottomrule
\end{tabular}
\end{center}

Total cost under 0.30\,EUR; no firmware required.

% ===========================================================================
\section{Pin Assignment}
\label{sec:pins}
% ===========================================================================

Edit \file{board-pin-mapping.pcf} and replace every pin number with the
actual pad number from your schematic before running nextpnr-ice40.

The \sig{CLK} signal must be assigned to a \textbf{Global Buffer (GB)
input} pin.  Valid GB pins on the iCE40UP5K SG48 package are:
\textbf{20, 28, 31, 38, 41, 43, 48}.  The placeholder in the current PCF
(pin~35) is \emph{not} a GB pin.

\textbf{Valid user I/O pin numbers for SG48} (39 total):
\begin{quote}\ttfamily
2, 3, 4, 6, 9, 10, 11, 12, 13, 14, 15, 16, 17, 18, 19, 20, 21,\\
23, 25, 26, 27, 28, 31, 32, 34, 35, 36, 37, 38, 39,\\
40, 41, 42, 43, 44, 45, 46, 47, 48
\end{quote}
Pins 1, 5, 7, 8, 22, 24, 29, 30, 33 are power/ground or reserved and
\textbf{must not} appear in the PCF file --- nextpnr rejects them with an
error.

\textbf{I/O budget:} the design uses 35 of the 39 user I/O pins; 4 are
free.

\begin{center}
\begin{tabular}{lr}
\toprule
\textbf{Group} & \textbf{Pins} \\
\midrule
Address A[7:0] (A15--A8 decoded by \sig{/IO1}, \sig{/IO2}) & 8 \\
Data D[7:0] (input only, shared net with SRAM DQ)          & 8 \\
\sig{PHI2}, \sig{RNW}                                      & 2 \\
\sig{IO1\_N}, \sig{IO2\_N}                                 & 2 \\
\sig{VERA\_CS}                                             & 1 \\
SRAM: \sig{SRAM\_PAGE[5:0]}, \sig{SRAM\_BLOCK[3:0]}, CE/OE/WE & 13 \\
\sig{CLK} (25\,MHz, GB pin)                                & 1 \\
\midrule
\textbf{Total}                                             & \textbf{35} \\
\bottomrule
\end{tabular}
\end{center}

\subsection{Complete pin table}

\begin{center}
\small
\begin{longtable}{lllllp{4.6cm}}
\toprule
\textbf{Signal} & \textbf{Dir} & \textbf{PCF pin} & \textbf{C64 port} & \textbf{Std} &
\textbf{Description} \\
\midrule
\endfirsthead
\toprule
\textbf{Signal} & \textbf{Dir} & \textbf{PCF pin} & \textbf{C64 port} & \textbf{Std} &
\textbf{Description} \\
\midrule
\endhead
\midrule
\multicolumn{6}{r}{\small continued on next page}\\
\endfoot
\bottomrule
\endlastfoot

% --- Clock ---
\sig{CLK}      & In  & 35* & --- & LVCMOS33 &
    25\,MHz system clock. Must land on a GB pin. \\[2pt]

% --- Address bus ---
\sig{A[0]}     & In  & 40* & Y & LVCMOS33 &
    Address bit 0. Also wired to SRAM A[0] (same for A1--A7). \\
\sig{A[1]}     & In  & 41* & X & LVCMOS33 & Address bit 1 \\
\sig{A[2]}     & In  & 42* & W & LVCMOS33 & Address bit 2 \\
\sig{A[3]}     & In  & 43* & V & LVCMOS33 & Address bit 3 \\
\sig{A[4]}     & In  & 44* & U & LVCMOS33 & Address bit 4 \\
\sig{A[5]}     & In  & 45* & T & LVCMOS33 & Address bit 5 \\
\sig{A[6]}     & In  & 46* & S & LVCMOS33 & Address bit 6 \\
\sig{A[7]}     & In  & 47* & R & LVCMOS33 & Address bit 7 \\[2pt]

% --- Data bus ---
\sig{D[0]}     & In  & 11* & 21 & LVCMOS33 &
    Data bit 0 (register writes). Shared net with SRAM DQ[0]
    (same for D1--D7). \\
\sig{D[1]}     & In  & 12* & 20 & LVCMOS33 & Data bit 1 \\
\sig{D[2]}     & In  & 13* & 19 & LVCMOS33 & Data bit 2 \\
\sig{D[3]}     & In  & 14* & 18 & LVCMOS33 & Data bit 3 \\
\sig{D[4]}     & In  & 15* & 17 & LVCMOS33 & Data bit 4 \\
\sig{D[5]}     & In  & 16* & 16 & LVCMOS33 & Data bit 5 \\
\sig{D[6]}     & In  & 17* & 15 & LVCMOS33 & Data bit 6 \\
\sig{D[7]}     & In  & 18* & 14 & LVCMOS33 & Data bit 7 \\[2pt]

% --- Control ---
\sig{PHI2}     & In  & 34* & E  & LVCMOS33 &
    Phase-2 clock. Synchronised internally; not used as FPGA clock. \\
\sig{RNW}      & In  & 36* & 5  & LVCMOS33 &
    Read/Not-Write. High = read, low = write. \\
\sig{IO1\_N}   & In  & 26* & 7  & LVCMOS33 &
    \sig{/IO1}: CPU access to \hex{DE00}--\hex{DEFF} (active low). \\
\sig{IO2\_N}   & In  & 27* & 10 & LVCMOS33 &
    \sig{/IO2}: CPU access to \hex{DF00}--\hex{DFFF} (active low). \\[2pt]

% --- VERA ---
\sig{VERA\_CS} & Out & 32* & --- & LVCMOS33 &
    VERA chip select (active low), per cycle: \sig{/IO2} low,
    \hex{DF00}--\hex{DF1F}, VERA selected (\hex{DFFD}), PHI2 high. \\[2pt]

% --- External SRAM ---
\sig{SRAM\_PAGE[0]}  & Out &  2* & --- & LVCMOS33 & SRAM A[8], page bit 0 \\
\sig{SRAM\_PAGE[1]}  & Out &  3* & --- & LVCMOS33 & SRAM A[9], page bit 1 \\
\sig{SRAM\_PAGE[2]}  & Out &  4* & --- & LVCMOS33 & SRAM A[10], page bit 2 \\
\sig{SRAM\_PAGE[3]}  & Out &  6* & --- & LVCMOS33 & SRAM A[11], page bit 3 \\
\sig{SRAM\_PAGE[4]}  & Out &  9* & --- & LVCMOS33 & SRAM A[12], page bit 4 \\
\sig{SRAM\_PAGE[5]}  & Out & 10* & --- & LVCMOS33 & SRAM A[13], page bit 5 \\
\sig{SRAM\_BLOCK[0]} & Out & 19* & --- & LVCMOS33 & SRAM A[14], block bit 0 \\
\sig{SRAM\_BLOCK[1]} & Out & 20* & --- & LVCMOS33 & SRAM A[15], block bit 1 \\
\sig{SRAM\_BLOCK[2]} & Out & 21* & --- & LVCMOS33 & SRAM A[16], block bit 2 \\
\sig{SRAM\_BLOCK[3]} & Out & 38* & --- & LVCMOS33 & SRAM A[17], block bit 3 \\
\sig{SRAM\_CE\_N}    & Out & 23* & --- & LVCMOS33 &
    SRAM chip enable (active low): \sig{/IO1} low and PHI2 high.
    Also the GeoRAM term of the data transceiver output enable. \\
\sig{SRAM\_OE\_N}    & Out & 39* & --- & LVCMOS33 &
    SRAM output enable (active low), read cycles only. \\
\sig{SRAM\_WE\_N}    & Out & 25* & --- & LVCMOS33 &
    SRAM write enable (active low), write cycles only. \\

\end{longtable}
\end{center}

\textit{* Pin numbers are placeholders.  Replace with actual pad numbers
from your board schematic.}

% ===========================================================================
\section{GeoRAM}
% ===========================================================================

\subsection{External SRAM wiring}

\textbf{AS6C4008-55} (Alliance Memory, 512K$\times$8, 55\,ns, 2.7--5.5\,V),
lower 256\,KB used.  The SRAM sits on the 3.3\,V side of the level
shifters, on the same nets as the FPGA, and is powered at 3.3\,V.

\begin{center}
\begin{tabular}{lll}
\toprule
\textbf{SRAM pin} & \textbf{Source} & \textbf{Note} \\
\midrule
A[7:0]   & C64 A[7:0] (PCB net)         & Direct trace, no FPGA I/O pin \\
A[13:8]  & FPGA \sig{SRAM\_PAGE[5:0]}   & Page register \hex{DFFE} \\
A[17:14] & FPGA \sig{SRAM\_BLOCK[3:0]}  & Block register \hex{DFFF} \\
A[18]    & GND (PCB)                    & Lower 256\,KB of the 512\,KB chip \\
DQ[7:0]  & C64 D[7:0] (PCB net)         & Shared with FPGA D[7:0] \\
CE\#     & FPGA \sig{SRAM\_CE\_N}       & \\
OE\#     & FPGA \sig{SRAM\_OE\_N}       & \\
WE\#     & FPGA \sig{SRAM\_WE\_N}       & \\
\bottomrule
\end{tabular}
\end{center}

\subsection{Address formula}

16 blocks of 16\,KB, each made of 64 pages of 256\,bytes:

\[
\text{offset} = \text{block} \times 16384 + \text{page} \times 256 + A[7:0]
\]

\begin{lstlisting}[style=vhdl,caption={GeoRAM decode in \file{C6502\_Interface.vhd}}]
sel_page   <= '1' when IO2_N='0' and A = x"FE" else '0';
sel_block  <= '1' when IO2_N='0' and A = x"FF" else '0';
sel_georam <= '1' when IO1_N='0' else '0';
\end{lstlisting}

None of these terms depends on the \hex{DFFD} latch: the GeoRAM is always
active.

\subsection{Data transceiver}

The board must drive the C64 data bus during GeoRAM reads.
\sig{SRAM\_CE\_N} is the GeoRAM term of the data transceiver output enable:

\[
\overline{OE} = \text{NAND}\bigl(\text{PHI2},\ \text{NAND}(\text{R/W},\
                \text{AND}(\sig{VERA\_CS},\ \sig{SRAM\_CE\_N}))\bigr)
\]

Every \hex{DExx} read is driven from the SRAM (GeoRAM always active);
\hex{DF00}--\hex{DF1F} reads are driven only while the VERA is selected.

% ===========================================================================
\section{Timing}
\label{sec:timing}
% ===========================================================================

The C64 runs at 0.985\,MHz (PAL) or 1.023\,MHz (NTSC); PHI2 is high for
$\approx$490\,ns.

\subsection{Register writes}

The 6510 drives write data during PHI2 high, up to $\approx$150\,ns after
the rising edge.  The latch (\hex{DFFD}) and the GeoRAM registers
(\hex{DFFE}/\hex{DFFF}) are sampled \cmd{SAMPLE\_CLKS} clock cycles after
the synchronised PHI2 rising edge: with the default of 6 clocks at
25\,MHz this is 240\,ns after \cmd{phi2\_rise}, i.e.\ $\approx$280--320\,ns
after the real edge, well inside the PHI2-high window.  A PHI2 pulse that
ends before the sample point never samples.

\subsection{SRAM read cycle}

\begin{center}
\begin{tabular}{lrl}
\toprule
\textbf{Event} & \textbf{Time from PHI2 rise} & \textbf{Note} \\
\midrule
\sig{/IO1} asserted by the PLA & $\approx$50--80\,ns & \\
\sig{SRAM\_CE\_N} asserted     & $\approx$50--80\,ns & \sig{/IO1} AND PHI2 \\
SRAM data valid (tAA = 55\,ns) & $\approx$105--135\,ns & AS6C4008-55 \\
6510 samples the data bus      & $\approx$490\,ns & At PHI2 fall \\
\bottomrule
\end{tabular}
\end{center}

The read margin is several hundred nanoseconds.

\subsection{SRAM write cycle}

\sig{SRAM\_WE\_N} is asserted combinationally while \sig{/IO1} is low,
PHI2 is high and R/W is low.  The pulse lasts $\approx$400\,ns, far above
the 40\,ns tWP of the AS6C4008-55.  The SRAM latches the data on the WE\#
rising edge (PHI2 fall); tDH = 0\,ns is covered by the 6510 data hold
after PHI2 fall.

\sig{SRAM\_PAGE}/\sig{SRAM\_BLOCK} change only on \hex{DFFE}/\hex{DFFF}
writes, which are \sig{/IO2} cycles, so they never change while CE\# is
low.

% ===========================================================================
\section{Source File Reference}
% ===========================================================================

\begin{center}
\begin{tabular}{lp{9cm}}
\toprule
\textbf{File} & \textbf{Description} \\
\midrule
\file{C6502\_Interface.vhd} &
    Top-level VHDL entity (Rev~6, C64 expansion port, VERA + GeoRAM).
    Contains the whole bus logic: \sig{/IO1}/\sig{/IO2} decode, VERA
    select latch, GeoRAM registers and SRAM control. \\
\file{board-pin-mapping.pcf} &
    Pin constraints for iCE40UP5K SG48.
    All pin numbers are placeholders -- edit before use. \\
\file{Makefile} &
    Targets: \cmd{all} (default), \cmd{flash}, \cmd{flash-iceprog},
    \cmd{doc}, \cmd{clean}. \\
\bottomrule
\end{tabular}
\end{center}

\end{document}
